本科毕业设计期间，首先要衷心感谢指导教师文静老师在课题方向、研究方法和论文撰写方面给予的耐心指导。老师严谨务实的科研态度和对细节的高标准要求，使我在项目实现和学术表达上都获得了显著提升。

感谢实验室同学在数据准备、环境配置和问题讨论中的帮助。每一次代码调试和实验复盘，都离不开大家的建议与支持。

感谢学院提供的计算资源与学习平台，使本课题能够顺利完成从模型训练到系统验证的全流程实践。

最后，感谢家人一直以来的理解与鼓励。你们的支持是我坚持完成本科阶段学习与毕业设计的动力来源。
